\documentclass[a4paper]{article}
\usepackage[utf8]{inputenc}
\usepackage[T1]{fontenc}

\title{}
\author{}
\date{}

\begin{document}


\section{Introduction}


Ces dernières 10aines d'années, les outils d'informatique graphique se sont fait
une place de choix dans notre quotidien. %
Ces outils, sont utilisés pour l'aménagement avec la modélisation d'espaces en
intérieur et en extérieur. Ils permettent de créer des objets virtuels, du
mobilier, mais aussi des personnages. %
Grâce à cela, nous pouvons, par exemple, regarder des films avec des effets
spéciaux remarquables et des films dont les personnages, plus ou moins
réalistes, sont animés. %


Autrement dit, ces outils, et plus particulièrement ceux dits de modélisation
géométrique, touchent de nombreux domaines d'applications. %
Parmi ces domaines, certains tels que l'archéologie et la conception assistée
par ordinateur (ou CAO) nécessite de créer des objets précis et de faire varier
leurs processus de constructions. %
Dans cet exemple, il s'agit d'ajouter une pergola sur une façade de maison, mais
pour les applications industrielles, il s'agit de concevoir toute une variété
de rouages.  %

Pour cela, nous nous intéressons plus particulièrement aux systèmes de
modélisation paramétrique. %
Ceux-ci présentent plusieurs avantages qui, à terme, permettent de faire varier
des constructions. %
Notamment, ces outils permettent de concevoir des pièces uniques à l'aide
d'opérations de modélisation géométriques. %
D'autre part, ces opérations, qui constituent alors ce que l'on appelle une
spécification paramétrique sont sauvegardés. %
Ainsi, il est possible de modifier les paramètres géométrique des opérations et
rejouer une construction. %


Modéliser un objet complexe, c'est un processus long où on enchaîne
plein d'opérations de modélisation.

Aujourd'hui, tous les modeleurs enregistrent ces processus dans ce qu'on appelle
des spécifications paramétriques pour les modifier et les rejouer. On peut
modifier les paramètres topologiques ou géométriques, voire ajouter, supprimer
ou déplacer des opérations.

On va voir un exemple de spécification qui sera notre fil rouge~:
\begin{itemize}
	\item Une première opération crée la face carrée f1 et ses arêtes dont e1.
	\item La suivante extrude f1 en un cube et e1 en la face f2.
	\item La triangulation de f2 crée plusieurs faces dont f7.
	\item La dernière opération colore f7.
\end{itemize}

Maintenant, on modifie cette spécification~: on ajoute l'opération qui
insère un sommet sur l'arête e1 et on la réévalue.

\begin{itemize}
	\item f1 et e1 sont recréées à l'identique
	\item Le sommet inséré scinde e1 en e5 et e6
	\item L'extrusion s'applique sur f1 mais cette fois e5 et e6 sont extrudés en f3 et
	      f4.
\end{itemize}

Et maintenant, comment on réévalue puisque f2 n'existe plus sous sa forme initiale,
pareil pour f7 ?

Intuitivement, f2 est scindée en f3 et f4 alors on voudrait trianguler f3 et/ou f4.

Autrement dit, mettre f2 comme référence, c'est pas robuste ici elle n'existe plus.

C'est là qu'on a besoin d'un nom persistant pour identifier les cellules par
leurs histoires et dire, par exemple, «~f3 et f4 sont construits comme f2, alors
on triangule ces deux faces~».



\section{Formalismes}



On commence avec les cartes généralisées, ou G-cartes, pour représenter des
objets



\subsection{G-cartes}


On prend un objet

Il se décompose en deux volumes connectés par des liaisons 3 en vert

Ses volumes se décomposent en faces connectées par des liaisons 2 en bleu

Qui se décomposent elles-mêmes en arêtes connectées par des liaisons 1 en rouge

Qui se décomposent en sommets connectés par des liaisons 0 en noir

Cette G-carte, c'est un graphe dont les nœuds sont des brins et les arcs des
liaisons.



\subsection{Orbites}



Dans les G-cartes, on représente les cellules avec des sous-graphes que l'on
appelle orbites.

Par exemple~:

Le sommet incident au brin «~a~» contient «~a~», tous les brins atteignables par des
liaisons 1, 2 et 3 et les liaisons elles-mêmes.

Cette arête contient le brin «~a~» et tous ceux atteignables par des liaisons 0,2,3.

Une face contient tous les brins atteignables par des liaisons 0,1,3

un volume contient ceux atteignables par des liaisons 0,1,2

Les orbites représentent aussi toutes sortes d'entités topologiques plus générales que des
cellules comme cette arête de volume.

Une orbite c'est un sous-graphe de type o incident à un brin ((b) en vert)



\subsection{Règles de transformation de graphe Jerboa}



Maintenant qu'on a formalisé les objets avec les G-cartes, on formalise les
opérations avec les règles de transformations de graphe Jerboa.

On prend l'exemple de la triangulation barycentrique d'une face.

De manière générale, une règle se compose d'une partie gauche qui filtre un
sous-graphe d'un objet à transformer et d'une partie droite qui décrit la transformation.

À gauche, le nœud n0 porte le type d'orbite face <0,1,3>.

Quand on applique la règle sur le brin a0, on filtre toute la face incidente à a0 en orange.

À droite, on retrouve le nœud n0, il est préservé donc ses brins sont préservés. La liaison 1 à gauche est remplacée par un «~underscore~» à droite, donc on supprime les liaisons 1 dans l'objet.

n2 est un nœud créé, donc ses brins sont créés comme des copies des brins filtrés par n0. %
Dans la règle, les liaisons 0 sont renommées en liaisons 1 et les liaisons 1 en liaisons 2. %
%Les liaisons 0, 1 de n0 à gauche sont renommés respectivement en 1,2 dans n2 à droite
%Donc, dans l'objet, il y a des liaisons 1 et 2 entre les brins.
Donc, dans l'objet, le sommet bleu et créé comme le dual de la face orange en remplaçant les liaisons 0,1 par des liaisons 1,2. %respectivement

Enfin, n1 est créé, avec des suppressions et du renommage.

Dans la règle, une liaison 1 connecte n0 à n1, ce qui lie les brins 2 à 2 dans l'objet.
Même principe pour la liaison 0 entre n1 et n2.

Les liaisons, dans les nœuds, sont appelées des arcs implicites.
Celles entre les nœuds sont des arcs explicites.

\subsection{Orbites}

Grâce aux règles, on peut suivre l'évolution des orbites.

On prend la face de volume incidente à n0.

À gauche, elle contient n0 avec ses arcs 0,1

À droite, elle contient n0 son arcs 0, l'arc 1 qui le connecte à
n1, n1, l'arc 0 qui connecte n1 à n2, n2 et son arc 1.

On remarque qu'il y a une seule orbite face dans la règle, alors qu'il y en a 4
dans l'objet. %

C'est parce que à gauche ce deuxième arc implicite 1, à droite, il est soit supprimé soit renommé en 2, en dehors du type d'orbite 0,1.
Donc on a une scission de la face filtrée. %

D'ailleurs, on voit qu'il n'y a pas de liaison 0,1 entre les faces



\section{Détection des évènements}



Venons-en au événements que l'on souhaite détecter à partir des formalismes
présentés plus tôt.

Nous allons vous présenter, dans cet ordre, comment repérer les événements de
créations, scission, fusion, suppression, modification et de non-changement à
partir des règles.

\subsection{Création}

Tout d'abord, une entité topologique est créée dans un objet lorsque n'existait
pas dans la version précédente de ce même objet.
Dans cet exemple, un sommet est créé sur une arête.

Ici, nous avons une règle où les nœuds n1 et n2 sont créés et filtrent donc les dimensions 0,1 et 2.
Suivant l'orbite R<02>(n1), cette règle crée au moins une arête.

Dans le graph H, n1 et n2 créent 4 arêtes dont les brins sont vert et bleu et connectés par des 0-arcs.

\subsection{Scission}

Une entité topologique est scindée dans un objet lorsque ses brins, entre G et
H, alors qu'ils sont préservés, cessent d'appartenir à une même cellule. Les
cellules résultantes de la scissions sont de la même nature.

Dans cette règle, nous présentons une scission explicite, c'est-à-dire qu'une
composante connexe est scindée en deux par la rupture d'un arc explicite.
Ici, la rupture d'un 2-arc sépare deux orbites <0>

G<02>(a0) est équivalent à G<02>(c1) et est composé de 4 brins.
H<02>(a0) n'en dispose plus que de 2 et est distinct de H<02>(c1).

Pour le cas implicite, L<01>(n0), autrement dit le nœud n0 a pour arcs
implicites 0 en première position et 1 en deuxième position. L'orbite associée
dans R comporte les nœuds n0, n1 et n2. Or aucun de ces nœuds n'a d'arc
implicite étiqueté 0 ou 1 en deuxième position. Il faut comprendre par là, qu'un
ensemble de brins est scindé en sous-ensembles distincts les uns des autres.

On peut voir que la face filtré G<01>(a0) a été scindée en 4 faces
H<01>(a0, c0, e0 et g0) toutes distinctes puisque 2 n'est pas suivi.

\subsection{Fusion}

Deux entités topologiques sont dites fusionnées lorsqu'au moins deux de leurs
brins, qui appartenaient à deux orbites distinctes appartiennent à une même
orbite.
Ici, nous avons la fusion de deux sommet par contraction d'une arête.

Cas explicite et symétrique à celui de la scission. Deux orbites <02> sont
reliées par couture d'un 2-arc entre elles.

De mêmes, deux arêtes G<02>(a0) et G<02>(c1) sont fusionnées en une arête
H<02>(a0).

Dans ce cas, implicite, l'orbite L<1>(n1) est composée des nœuds n0 et n1. Aucun
de ces nœuds n'a d'1-arc implicite. Or, R<1>(n1) a pour deuxième arc implicite
un 1-arc. Il y a donc fusion.

Dans cet exemple, nous mettons en évidence la fusion de sommets de faces. En
effet, les sommets G<1>(j1) et G<1>(d1) sont fusionnés en H<1>(j1) et les
sommets G<1>(m1) et G<1>(c1) sont fusionnés en H<1>(m1).

Notons un cas, limite. Ici, deux faces sont fusionnées, cependant L<01>(n1) est
incomplète, la fusion ne peut donc être garantie et nécessite une confirmation
en analysant les cellules correspondantes.

\subsection{Autres événements}

Nous abordons désormais les trois derniers événements

\subsection{Suppression}

La suppression est en quelque sorte le symétrique de la création, pour qu'une orbite soit entièrement supprimée, il est nécessaire qu'elle soit complète.

Ici, l'orbite sommet L<02>(n0) est supprimée. Il en résulte, dans G la
suppression de l'arête d'orbite G<02>(a0) dont tous les brins sont supprimés
dans H.

\subsection{Modification}

La modification est un événement où au moins un arc d'une orbite est réécrite,
avec suppression ou création de nœud, et ne résulte ni en scission ni en fusion.

Ici, l'orbite sommet L<1,2>(n1) est modifiée car il existe un deuxième arc
implicite écrit 2 dans n0 et dans R<1,2>(n1), le deuxième arc implicite de n1
est réécrit en 2. Il n'y a donc ni scission, ni fusion et il y a réécriture
d'un arc implicite avec une étiquette issue du type <1,2>.

\subsection{Non-changement}

Cet événement est le symétrique de la modification. Une orbite est inchangée si
elle ne subit aucune réécriture sur son type et qu'elle n'est ni augmentée ou
réduite.

Ici, L<0>(n0) filtre des arêtes de faces et puisque dans R<0>(n0) le premier arc
implicite de n0 est toujours 0, l'orbite est inchangée.

Concrètement, les arêtes de la face triangulée sont inchangées : relation
identiques entre les mêmes brins.

\section{Rejeu à base de règles}


\subsection{Nomination Persistante}

Avec ces formalismes, on peut maintenant voir la nomination persistante.
On reprend l'exemple de l'introduction.



% \subsubsection{Nom persistant brin}



Ici, on propose d'enregistrer l'histoire des brins comme étant leurs noms
persistants. %
%Chaque brin contient alors les numéros d'opérations et de nœuds qui les créent
%ou réécrivent.

Cette première règle crée une face carrée.
Elle a 8 nœuds dont n3 et n6.

Dans cette première application, n3 crée le brin 3, son histoire c'est le numéro de l'application suivi du nom du nœud qui le crée, elle se lit [1n3] et on l'enregistre dans le brin 3.

n6 crée le brin 6, on enregistre son histoire [1n6].

Quand on applique l'extrusion sur le carré, on sélectionne le brin 6 pour désigner sa face incidente. %
Sauf qu'à la réévaluation le brin pourrait changer, alors dans la spécification, à cet instant, on enregistre une photographie de son nom persistant [1n6].

n1 à gauche filtre les brins 3 et 6 et n2 et n4 à droite créent les copies de 3, 33 et 35. %
On construit leurs histoires à partir de celle du brin 3. %

L'histoire de 33 c'est [1n3;2n2]. %
Celle de 35 c'est [1n3;2n4]. %
L'histoire de 3 devient [1n3;2n1], pareil pour le brin 6. %

Ensuite, pour trianguler la face de devant, on sélectionne le brin 33 pour désigner sa face incidente et on enregistre son nom persistant.

À droite, n0 réécrit les brins 33 et 35.
C'est la troisième application, l'histoire du brin 33 devient [1n3;2n2;3n0], même principe pour 35.% qui devient [1n3;2n4;3n0].

Enfin, on colorie la face incidente au brin 35 et on enregistre son nom persistant.

On vient de nommer des brins de manière unique en représentant leurs histoires,
il s'agit de la nomination persistante des brins.

Mais, dans une spécification, les paramètres topologiques sont des orbites.
Donc, on doit représenter les histoires de ces orbites.



% \subsubsection{Nom persistant orbite}



Pour reconstituer l'histoire d'une orbite, on remonte
l'histoire du nom persistant de brin qui la désigne.

On va prendre l'orbite désignée par [1n3;2n4;3n0] dans l'opération de coloration.

Les arcs implicites de ce n0, en orange, désignent une face de volume.

On commence l'histoire par la fin à la troisième opération appliquée qui est la triangulation.
On s'intéresse à la face de volume incidente à n0 ce qui correspond à la face orange dans la règle

Comme on l'a vu plus tôt, cette face en orange provient de la scission de la face filtrée à gauche et a pour origine une arête.

On reporte ces informations dans l'histoire de l'orbite.

Donc la face qui nous intéresse, c'est la face incidente à n0, en orange.
Elle est produite par la triangulation, qui scinde la face filtrée à gauche, notée avec la flèche noire, et a pour origine une arête notée par la flèche rouge.

On remonte l'histoire d'un cran à la deuxième opération, qui est l'extrusion, et son nœud n4.

Ce qu'on doit suivre maintenant, c'est la face et l'arête incidentes à n4, notées par les boîtes bleue et orange dans la règle.

La face est créée par l'extrusion d'une arête et l'arête par l'extrusion d'un brin.

Enfin, on remonte encore l'histoire, cette fois à la toute première opération de création de carré et son nœud n3.

Les orbites que l'on suit sont créées par la règle de création de carré.

On vient de reconstituer l'histoire d'une orbite à partir du nom persistant d'un brin et des règles de la spécification.

On fait ça uniquement à partir de la spécification et de ses règles, sans les réappliquer et sans utiliser l'objet.

On suit la même procédure pour chacun des noms persistants de brins enregistrés comme paramètres topologiques enregistrés dans la spécification et uniquement ceux-ci.

%Donc les histoires des orbites, on les reconstitue à la demande seulement si elles sont utilisées dans la spécification.

Donc on ne reconstitue l'histoire des orbites que si elles sont utilisées dans la spécification.

Avec les noms persistants d'orbites, on peut maintenant passer à la réévaluation.



% \section{Réévaluation}



On reprend l'histoire précédente, cette fois de haut en bas, et on réévalue
la spécification modifiée. %

La première opération recrée la face carrée à l'identique.

Le nœud n3 de la règle de création de carré crée un unique brin, 3.

L'opération suivante a été ajoutée, j'affiche pas la règle mais si on l'analyse on voit qu'elle scinde en deux l'arête filtrée.
À partir de ses nœuds on retrouve les deux arêtes scindées, l'une désignée par le brin 3, l'autre par le brin 4.

Donc on a une histoire par arête. %

Pour l'opération d'extrusion, on aura déterminé son paramètre topologique au
préalable dans un autre arbre, exactement comme on fait là.

Chaque arête est extrudée en une nouvelle face, et chaque brin en une
arête.

Ici, n4 crée plusieurs brins, mais une seule copie de 3, qui est 37, et une seule
copie de 4, qui est 52.



% \subsection{stratégies de rejeu}



Enfin, on peut réappliquer la triangulation sur chaque branche.

n0 réécrit 37 et 52 à l'identique.

C'est la dernière application enregistrée, 37 et 52 désignent les faces à colorer.

Notre procédure permet de retrouver tous les paramètres potentiels, ici on peut laisser le choix à l'utilisateur de réappliquer les opérations sur l'une et/ou l'autre branche. %

\section{Extension aux scripts}

\section{Conclusion}


\end{document}
